\chapter{Introdução}
\label{cap:introducao}

O mapeamento tridimensional de ambientes confinados é relevante para inspeção, documentação e exploração de locais onde a presença humana pode ser arriscada ou operacionalmente custosa. Cavernas, galerias, áreas industriais internas e regiões de difícil acesso têm em comum a baixa disponibilidade de infraestrutura externa de posicionamento. Em particular, sistemas baseados em GPS tendem a apresentar indisponibilidade ou degradação severa nesses cenários, exigindo métodos de localização relativa baseados em sensores embarcados \cite{thrun2005probabilistic,beardMcLain2012}.

Uma solução direta seria integrar as acelerações medidas por uma Unidade de Medição Inercial (IMU) para obter velocidade e posição. Entretanto, essa abordagem isolada é limitada: pequenos vieses, ruídos de quantização e erros de atitude são integrados ao longo do tempo, produzindo deriva crescente na estimativa de posição \cite{tittertonWeston2004,groves2013principles,farrell2008aided}. Por outro lado, sensores ópticos de fluxo medem deslocamentos relativos da superfície observada entre quadros consecutivos, fornecendo uma fonte direta de odometria no plano observado, especialmente útil em navegação indoor e em ambientes sem GPS \cite{lucasKanade1981,pixart2017pmw3901}.

Este trabalho adota uma arquitetura híbrida. O sensor óptico de fluxo é tratado como fonte primária de deslocamento no plano horizontal, com correção de escala pela altura medida por um sensor ToF. A IMU complementa essa estimativa ao fornecer atitude do drone e ao preencher pequenos intervalos nos quais o sensor óptico não registra evento de movimento. O LiDAR 2D, por sua vez, fornece distância, ângulo e intensidade de retorno para cada ponto de varredura. Com a posição e a atitude estimadas, cada ponto LiDAR é transformado do referencial local do sensor para o referencial global, resultando em uma nuvem de pontos tridimensional.

\section{Problema de pesquisa}

O problema investigado pode ser resumido na seguinte pergunta: como estimar a posição relativa de uma plataforma móvel em ambiente sem GPS e usar essa pose para transformar medições LiDAR em uma nuvem de pontos global coerente? A resposta proposta combina três fontes de informação com papéis distintos: fluxo óptico para deslocamento horizontal, IMU para orientação e propagação em intervalos sem medição óptica, e LiDAR para geometria do ambiente.

\section{Objetivos}

O objetivo geral deste trabalho é desenvolver e documentar um fluxo de aquisição e processamento capaz de estimar a localização relativa de um drone em ambiente sem GPS e gerar uma nuvem de pontos tridimensional no referencial global a partir da combinação de sensor óptico, IMU e LiDAR.

Os objetivos específicos são:

\begin{itemize}
    \item Integrar em uma plataforma embarcada sensores de fluxo óptico, ToF, IMU e LiDAR, registrando os dados em formato sincronizado por timestamp.
    \item Converter e filtrar os sinais inerciais para estimar atitude e aceleração linear no referencial global.
    \item Estimar a trajetória no plano $XY$ usando fluxo óptico como fonte primária e integração inercial apenas nos intervalos sem eventos do sensor óptico.
    \item Processar os pontos LiDAR com filtros físicos e estatísticos, convertendo distância e ângulo para coordenadas locais.
    \item Transformar a nuvem local do LiDAR para o referencial global usando posição, atitude e hipóteses de montagem do sensor.
    \item Identificar limitações experimentais e parâmetros que ainda dependem de calibração final.
\end{itemize}

\section{Contribuições}

As principais contribuições deste trabalho são:

\begin{itemize}
    \item Definição de uma arquitetura de sensoriamento para navegação relativa em ambiente sem GPS, usando fluxo óptico como odometria primária e IMU como apoio inercial.
    \item Definição de um fluxo de processamento para conversão de unidades, calibração de viés do giroscópio, filtragem Butterworth, estimação de atitude por Madgwick e fusão de posição no plano $XY$.
    \item Definição de um fluxo de geração de nuvem global que considera a rotação do LiDAR principal em torno do eixo $Y$ do drone e aplica filtros automáticos de intensidade, distância e continuidade de pose.
    \item Organização das hipóteses e limitações do experimento, especialmente a ausência de odometria vertical dedicada e a necessidade de validação da calibração extrínseca entre sensores.
\end{itemize}

\section{Organização do trabalho}

O Capítulo~\ref{cap:conceitos} apresenta os conceitos de sensores embarcados, navegação inercial, fluxo óptico, atitude, transformações de coordenadas e filtragem usados no projeto. O Capítulo~\ref{cap:metodologia} descreve a plataforma, os dados coletados e o método de processamento adotado. O Capítulo~\ref{cap:resultados_conclusao} apresenta a implementação, os resultados obtidos, as limitações e os trabalhos futuros.
